\scanpage{168}
\mspage{26}

est en fait cofiltrante\deleted{\unsure{, donc}}, \unsure{Donc|donc} pour deux objets $a,b$ de $S$, $\exists\ c$ avec $c\longrightarrow a$, $c\unsure{\mathbin{-}|\longrightarrow}b$ (condition PS1$'$, plus forte que PS1). Pour prouver que $\unsure{\mathscr X|\mathcal X}$ est tot. $0$-connexe, on est réduit à prouver \inserted{par 2.3.} que les $a\times b$ dans $\mathscr X^\wedge$ sont $0$-connexes. Par 2.9, il revient au $=$ de prouver qu'ils sont $0$-connexes dans $S^\wedge$. Donc on est ramené à prouver le

\uline{Lemme 2.10}. Soit $S$ une catégorie cofiltrante. Si $a,b$ sont deux objets de $S$, $a\times b$ est un objet $0$-connexe de $S^\wedge$.

On voit de suite qu'il \deleted{\unsure{revient au même de}} \inserted{suffit de prouver (et} en fait il revient au $=$, compte tenu que $a\times b$ est loc. \deleted{\unsure{donc}} \inserted{irréductible)} que si on a deux flèches

% slanted
\margin{\deleted{\uline{\unsure{Lemme}} Soit $f$ \unsure{un|une} \illegible{1}. Alors $f$ \illegible{2} \unsure{est propre à}}}

\diagram{\begin{tikzcd}[ampersand replacement=\&]
x \arrow[r] \& a\times b \\
\& y \arrow[u]
\end{tikzcd}}

alors $x\times_{a\times b}y\ne\varnothing$ i.e $\exists\,z$ dans $S$ rendant commutatif

\diagram{\begin{tikzcd}[ampersand replacement=\&]
x \arrow[r] \& a\times b \\
z \arrow[u] \arrow[r] \& y \arrow[u]
\end{tikzcd}}

Or par PS1$'$ il existe bien un $z$ s'envoyant dans $x$ et dans $y$, mais les deux flèches composées

\[
z\rightrightarrows a\times b
\]

i.e. donnant des doubles flèches

\[
z\rightrightarrows a \tag*{$(*)$}
\]

\[
z\rightrightarrows b, \tag*{$(*')$}
\]

peuvent n'être pas égales. Mais par PS2,

% archive sheet number: 168
